\section{Boolean-Kernel Evaluation}
\label{sec:evaluation}

We now derive the single-point evaluation algorithm for the
Boolean-kernel representation introduced in \cref{sec:representations}.

Let
\[
P(X)
=
\sum_{y\in\{0,1\}^m}
c_y
\prod_{i=0}^{m-1}
\left(X^{2^i}+y_i\right),
\]
and fix an evaluation point
\[
x\in\mathbb{F}.
\]
Define
\[
t_i=x^{2^i},
\qquad
0\le i<m.
\]

The key observation is that two basis polynomials whose indices differ
only in one Boolean coordinate can be combined locally.

\subsection{The elementary fold identity}

Fix a level \(i\).  Let
\[
u=(u_0,\ldots,u_{i-1},u_{i+1},\ldots,u_{m-1})
\in\{0,1\}^{m-1}
\]
denote the common values of all Boolean coordinates other than \(i\).
Let
\[
K_{u,0}(X)
\]
and
\[
K_{u,1}(X)
\]
denote the two kernel polynomials whose \(i\)-th coordinate is respectively
\(0\) and \(1\).

By definition,
\[
K_{u,0}(X)
=
X^{2^i}
\prod_{j\ne i}
\left(X^{2^j}+u_j\right),
\]
while
\[
K_{u,1}(X)
=
\left(X^{2^i}+1\right)
\prod_{j\ne i}
\left(X^{2^j}+u_j\right).
\]

Suppose the corresponding coefficients are
\[
a,b\in\mathbb{F}.
\]
Their contribution to \(P(x)\) is
\[
aK_{u,0}(x)+bK_{u,1}(x).
\]
Factoring the common part gives
\[
\begin{aligned}
aK_{u,0}(x)+bK_{u,1}(x)
&=
\left(
a\,x^{2^i}
+
b(x^{2^i}+1)
\right)
\prod_{j\ne i}
\left(x^{2^j}+u_j\right)
\\
&=
\left(
x^{2^i}(a+b)+b
\right)
\prod_{j\ne i}
\left(x^{2^j}+u_j\right).
\end{aligned}
\]

Therefore, after eliminating Boolean coordinate \(i\), the two
coefficients \(a\) and \(b\) are replaced by the single coefficient
\[
\boxed{
F_{t_i}(a,b)
=
t_i(a+b)+b
}
\]
with
\[
t_i=x^{2^i}.
\]

\begin{lemma}[Pair-fold identity]
\label{lem:pair-fold}
Let \(a,b\in\mathbb{F}\), let \(i\in\{0,\ldots,m-1\}\), and let \(u\)
fix all Boolean coordinates other than \(i\).  Then
\[
aK_{u,0}(x)+bK_{u,1}(x)
=
F_{t_i}(a,b)
\prod_{j\ne i}
\left(x^{2^j}+u_j\right),
\]
where
\[
F_{t_i}(a,b)=t_i(a+b)+b.
\]
\end{lemma}

\begin{proof}
The identity follows directly from the factorization above.
\end{proof}

\subsection{Recursive reduction}

The pair-fold identity can be applied recursively over all Boolean
coordinates.  We use the coefficient ordering in which the integer index
\[
y=\sum_{i=0}^{m-1}y_i2^i
\]
stores the Boolean vector
\[
(y_0,\ldots,y_{m-1}).
\]
Hence the least significant bit corresponds to the first folding level.

Let
\[
v^{(0)}_j=c_j,
\qquad
0\le j<N.
\]
After level \(i\), the active vector has length
\[
N_i=\frac{N}{2^i}.
\]

For
\[
0\le j<\frac{N_i}{2},
\]
define
\[
v^{(i+1)}_j
=
F_{t_i}
\left(
v^{(i)}_{2j},
v^{(i)}_{2j+1}
\right).
\]
Equivalently,
\[
v^{(i+1)}_j
=
t_i
\left(
v^{(i)}_{2j}
+
v^{(i)}_{2j+1}
\right)
+
v^{(i)}_{2j+1}.
\]

After \(m\) levels only one field element remains.

\begin{theorem}[Correctness of Boolean-kernel evaluation]
\label{thm:kernel-eval-correct}
Let
\[
P(X)
=
\sum_{y\in\{0,1\}^m}
c_yK_y(X).
\]
Construct the vectors
\[
v^{(0)},v^{(1)},\ldots,v^{(m)}
\]
by the recurrence above using
\[
t_i=x^{2^i}.
\]
Then
\[
v^{(m)}_0=P(x).
\]
\end{theorem}

\begin{proof}
We prove a stronger invariant.

After completing levels
\[
0,1,\ldots,i-1,
\]
the vector \(v^{(i)}\) is indexed by the remaining Boolean coordinates
\[
(y_i,\ldots,y_{m-1}).
\]
For every remaining index
\[
z=(z_i,\ldots,z_{m-1})\in\{0,1\}^{m-i},
\]
the value \(v^{(i)}_z\) is the coefficient of
\[
\prod_{j=i}^{m-1}
\left(X^{2^j}+z_j\right)
\]
after the already-eliminated coordinates
\[
0,\ldots,i-1
\]
have been evaluated at \(x\).

The invariant is true for \(i=0\) by the original representation of
\(P\).

Assume it holds for some \(i<m\).  Pair entries whose remaining indices
agree in every coordinate except \(y_i\).  By
\cref{lem:pair-fold}, replacing such a pair
\[
(a,b)
\]
by
\[
F_{t_i}(a,b)
\]
exactly evaluates the factor associated with coordinate \(i\), while
leaving the factors for coordinates
\[
i+1,\ldots,m-1
\]
unchanged.  Therefore the invariant holds for \(i+1\).

After \(m\) levels no Boolean coordinates remain, so the unique surviving
value is exactly
\[
P(x).
\]
\end{proof}

\subsection{Exact arithmetic complexity}

At level \(i\), the active vector has length
\[
\frac{N}{2^i},
\]
so the number of folds is
\[
\frac{N}{2^{i+1}}.
\]
Summing over all \(m\) levels gives
\[
\sum_{i=0}^{m-1}
\frac{N}{2^{i+1}}
=
N
\sum_{k=1}^{m}
2^{-k}
=
N\left(1-\frac{1}{2^m}\right)
=
N-1.
\]

Therefore the evaluation tree contains exactly
\[
\boxed{N-1}
\]
folds.

Each fold
\[
F_t(a,b)=t(a+b)+b
\]
contains, at the field-operation level,

\[
1
\]
multiplication and
\[
2
\]
additions.

Hence the tree itself requires exactly
\[
\boxed{N-1}
\]
field multiplications and
\[
\boxed{2(N-1)}
\]
field additions.

The level parameters are generated by
\[
t_0=x,
\qquad
t_{i+1}=t_i^2.
\]
Thus producing all
\[
t_0,\ldots,t_{m-1}
\]
requires
\[
\boxed{m-1}
\]
field squarings.

We summarize the operation count in the following proposition.

\begin{proposition}[Exact operation count]
\label{prop:operation-count}
For \(N=2^m\), Boolean-kernel single-point evaluation requires
\[
N-1
\]
kernel folds.  Using the direct fold formula, this corresponds to
\[
N-1
\]
field multiplications,
\[
2(N-1)
\]
field additions, and
\[
m-1
\]
field squarings for generation of the level parameters \(t_i\).
\end{proposition}

\subsection{Parallel structure}

The dependency graph is a complete binary reduction tree.  At each fixed
level \(i\), all folds operate on disjoint coefficient pairs and are
therefore independent.

The number of available independent folds at level \(i\) is
\[
\frac{N}{2^{i+1}}.
\]
Thus the lower levels expose abundant data parallelism, while the final
levels contain progressively less work.

A naive parallel implementation may synchronize globally after every tree
level.  For large \(N\), however, such synchronization can dominate the
shrinking amount of useful work near the top of the tree.  The optimized
implementation studied later therefore groups multiple consecutive levels
inside independent power-of-two subtrees and parallelizes only across
those larger subtrees.

This changes the scheduling strategy but not the arithmetic graph or the
number of folds.

\subsection{Input complexity}

The evaluation algorithm uses
\[
O(N)
\]
field operations and reads an arbitrary vector of \(N\) coefficients.

For the unrestricted problem considered here, this linear dependence on
the input length is unavoidable in the worst case.  To see this, suppose
an algorithm does not inspect some coefficient \(c_y\).  For any evaluation
point \(x\) such that
\[
K_y(x)\ne0,
\]
changing only \(c_y\) changes
\[
P(x)
\]
while leaving all inspected coefficients unchanged.  Hence an algorithm
that is correct on arbitrary coefficient vectors must in general obtain
information about every coefficient.

Accordingly, the principal optimization target is not the asymptotic
complexity but the constant factor: field reduction, memory traffic,
allocation, bounds checking, scheduling, and the organization of the fold
tree.

\subsection{Example for \(N=8\)}

Let
\[
m=3,
\qquad
N=8.
\]
The input coefficient vector is
\[
(c_0,c_1,c_2,c_3,c_4,c_5,c_6,c_7).
\]
Set
\[
t_0=x,
\qquad
t_1=x^2,
\qquad
t_2=x^4.
\]

The first level computes
\[
\begin{aligned}
u_0&=F_{t_0}(c_0,c_1),\\
u_1&=F_{t_0}(c_2,c_3),\\
u_2&=F_{t_0}(c_4,c_5),\\
u_3&=F_{t_0}(c_6,c_7).
\end{aligned}
\]

The second level computes
\[
\begin{aligned}
w_0&=F_{t_1}(u_0,u_1),\\
w_1&=F_{t_1}(u_2,u_3).
\end{aligned}
\]

Finally,
\[
P(x)=F_{t_2}(w_0,w_1).
\]

The tree uses
\[
4+2+1=7=N-1
\]
folds, illustrating the general count of
\cref{prop:operation-count}.
